\section{Selected reverse programs and sufficient traces}
\label{app:theory}

This appendix states the mathematical contract independently of a particular
framework. All vector spaces are finite-dimensional real spaces; arithmetic
equalities below are real-arithmetic identities. A floating-point
implementation must separately validate tie selection, exceptional values,
reduction order, and numerical error. In particular, equality of the
mathematical reverse actions does not imply bitwise equality after changing
the order of floating-point accumulation.

\subsection{Local contract and graph composition}
\label{app:composition}

Let an executed computation be a finite directed acyclic multigraph.
An incoming edge denotes an argument occurrence, so the same predecessor may
appear twice in one operation. At node $v$, a declared local differentiation
rule specifies linear maps $A_{ve}$ from the state on incoming edge $e$ to
the output state, and $B_v$ from the local parameter vector to that output.
These maps are evaluated at the recorded forward execution. They need not be
classical derivatives at nondifferentiable points; the declaration of the
rule is part of the contract.

Write $S$ for information available to reverse execution independently of the
certificate: graph and address metadata, shapes, any required parameter
snapshot, and explicitly retained boundary data. Local parameters may be
coordinate views of a shared global vector $\theta$, with
$\theta_v=E_v\theta$ for a fixed coordinate-selection map $E_v$.

\begin{definition}[Joint trace sufficiency]
\label{def:joint-trace}
A local record $c_v$ is sufficient for a requested local reverse action when
evaluators using only $(q,c_v,S)$ return $A_{ve}^{T}q$ for every requested
state edge $e$ and $B_v^{T}q$ for the requested local parameters, for every
incoming cotangent $q$. Every additional realized forward quantity used by an
evaluator must belong to either $c_v$ or the declared boundary information
$S$. A record that determines state actions need not determine parameter
actions.
\end{definition}

\begin{theorem}[Composition of trace-sufficient local reverse actions]
\label{thm:appendix-composition}
Suppose all local state actions needed by a finite executed graph, and all
requested local parameter actions, satisfy Definition~\ref{def:joint-trace}.
Given output cotangents, the records $(c_v)_v$, and $S$, reverse-topological
accumulation reproduces the declared local reverse program without access to
any unrecorded primal activation or forward recomputation. Shared parameters,
repeated argument occurrences, and fan-out are handled by summing their
contributions. If all local maps are classical derivatives at the execution
point, the result is the classical vector--Jacobian product of the graph.
\end{theorem}

\begin{proof}
Initialize every state and global-parameter cotangent to zero, then add the
supplied cotangents at the requested graph outputs. Process nodes in reverse
topological order. When $v$ is visited, every consumer of its value has
already been visited, so $\bar z_v$ contains exactly the accumulated
contributions from all downstream occurrences and any supplied output seed.
For every incoming argument occurrence $e$ with predecessor $u(e)$, perform
\begin{equation}
 \bar z_{u(e)}\mathrel{+}=A_{ve}^{T}\bar z_v,
 \qquad
 \bar\theta\mathrel{+}=E_v^{T}B_v^{T}\bar z_v
 \label{eq:reverse-accumulation}
\end{equation}
where the parameter update is performed once per node, not once per state
edge. The assumed evaluators provide precisely these terms from the record
and available information. Induction over reverse order establishes equality
with the declared reverse program. Repeated edges add repeated terms, and
$E_v^{T}$ scatters each local contribution into the correct shared parameter
coordinates. When the local maps are classical derivatives, the ordinary
chain rule identifies this program with the graph VJP.
\end{proof}

This is a specialization of standard reverse accumulation, not a new chain
rule or learning rule. The substantive operator question is whether a small
record determines both local state and local parameter actions. The global
selected linear map need not be a partial permutation: residual addition
duplicates cotangents, and overlapping routes can produce integer
multiplicities greater than one.

\begin{corollary}[Finite recurrent windows]
\label{cor:finite-window}
The theorem applies to any stated finite recurrent unrolling whose executed
nodes satisfy the local contract. It includes the cotangent of the initial
state when that state is differentiated, and sums parameter contributions
across all occurrences of a shared recurrent parameter.
\end{corollary}

\begin{proof}
Replace each recurrent invocation in the finite window by its time-indexed
node. The resulting graph is finite and acyclic; apply
Theorem~\ref{thm:appendix-composition}.
\end{proof}

The corollary makes no infinite-horizon, fixed-point, or convergence claim.
A bounded tape ring permits truncated windows, but overwriting earlier records
does not preserve an unrestricted full-history gradient. The current
\texttt{morphssm\_tape} adapter rejects recurrent scans and aliased parameters;
the corollary and alias construction describe mathematical composition, not
additional adapter coverage.

\subsection{Additive Extrema, Signs, and Top-\texorpdfstring{$k$}{k} Sums}
\label{app:operators}

For an independent example index $a$, consider
\begin{equation}
 y_{aj}=\mathop{\mathrm{ext}}_{i=1}^{K}
          (x_{ai}+\sigma w_{ji})+b_j,
 \qquad \mathrm{ext}\in\{\max,\min\},\quad \sigma\in\{-1,+1\}.
 \label{eq:signed-additive-extreme}
\end{equation}
An exact tie selects the lowest index in a fixed candidate ordering. Let
$r_{aj}$ be the selected index and let $q_{aj}$ be the output cotangent.

\begin{proposition}[Joint route sufficiency for additive extrema]
\label{prop:additive-route}
Under the selected-branch rule, the indices $r_{aj}$ and static addresses
determine all state, weight, and bias reverse actions of
\eqref{eq:signed-additive-extreme}:
\begin{align}
 \bar x_{ai}&=\sum_j q_{aj}\mathbf{1}\{r_{aj}=i\},\\
 \bar w_{ji}&=\sigma\sum_a q_{aj}\mathbf{1}\{r_{aj}=i\},\\
 \bar b_j&=\sum_a q_{aj}.
 \label{eq:additive-reverse}
\end{align}
No candidate activation value or weight value is required for these actions.
\end{proposition}

\begin{proof}
After selecting $r_{aj}$, the local branch is
$x_{a,r_{aj}}+\sigma w_{j,r_{aj}}+b_j$. Its nonzero derivatives are $1$,
$\sigma$, and $1$, respectively. Transposing and accumulating across outputs
and parameter-sharing examples gives the displayed formulas.
\end{proof}

Changing max to min changes the comparison, not the derivative sign. A minus
sign arises only from a minus sign in the selected affine candidate, as in
the common erosion convention $x-w$. TikiTorch's additive min-plus reference
uses $x+w$, so its selected weight derivative has positive sign.
Morphological convolution replaces $i$ by a static channel/offset address.
Overlapping windows and shared kernel coordinates scatter-add their
contributions. A selected padding constant receives no requested input
cotangent, although a selected additive kernel parameter still does.

For the sum of the $k$ selected candidates in each output, the same formula
replaces the single-index indicator by membership in the selected set. Its
coefficients are $1$, not $1/k$; averaging would be a different operator.
Stable ordered indices implement a deterministic tie contract, but ordering
within a selected set is unnecessary to determine the real-arithmetic VJP of
the sum. This distinction matters to storage lower bounds.

For $h_t=\max(u_t,h_{t-1}-\lambda)$, a carry/input gate determines the state
transpose action; each carry contributes the negative accumulated cotangent
to $\lambda$. If a different parameterization uses $|\lambda|$, its declared
sign selection must also be retained or otherwise available. A route index
does not automatically encode such signs.

\subsection{Ties: program agreement is not a global subgradient theorem}
\label{app:ties}

At a strict max/min winner, the selected branch is locally the active branch.
At a tie, each individual affine branch gradient is a valid local generalized
gradient of that extremum. Arbitrary composition of those selections need
not belong to the Clarke generalized Jacobian of the composite function.
The distinction between differentiating a program and differentiating its
denoted function is central in nonsmooth AD \citep{bolte2020model}.

\paragraph{An incompatible-tie counterexample.}
Let
\begin{equation}
 f(x)=\max(x,0)-\max(-x,0)=x.
\end{equation}
Order the candidate lists as $(x,0)$ and $(-x,0)$. At $x=0$, the first-index
rule selects the variable branch in both maxima. The composed local reverse
program returns $1-(-1)=2$, whereas the composite function has derivative
$f'(0)=1$ and Clarke generalized gradient $\{1\}$. Both selected branches
cannot be strict winners at a common nearby nonzero input. Deterministic
local ties therefore suffice for reproducible program replay, but do not
prove that the replay is a generalized gradient of the composite function.

A classical-VJP claim is valid when all relevant local functions are
differentiable at the recorded execution. For a continuous piecewise-affine
program, another sufficient condition at a boundary is that the complete
selected branch is realized on a sequence of full-dimensional affine regions
approaching the point: its constant Jacobian is then a limiting Jacobian of
the composite. We do not assert this simultaneous realizability for arbitrary
ties. Tests at ties must compare the explicitly declared local program;
ordinary central finite differences are not an oracle for that declaration.

\subsection{A local route-stability certificate}
\label{app:margin}

\begin{proposition}[Sufficient score-gap condition]
\label{prop:score-gap}
Let $r$ be the unique maximum of finite scores $(s_i)_i$. Suppose perturbed
scores obey $|s'_i-s_i|\leq\epsilon_i$. If
\begin{equation}
 s_r-s_i>\epsilon_r+\epsilon_i \qquad\text{for every }i\ne r,
\end{equation}
then $r$ is still the unique maximum. In particular, a winner--runner-up gap
$\gamma>2\epsilon$ suffices when every $\epsilon_i\leq\epsilon$. Reversing
signs gives the analogous statement for a minimum. For top-$k$ sums, a gap
greater than $2\epsilon$ between ranks $k$ and $k+1$ preserves the selected
set, but need not preserve its internal rank ordering.
\end{proposition}

\begin{proof}
For every competitor,
$s'_r-s'_i\geq s_r-s_i-\epsilon_r-\epsilon_i>0$.
For top-$k$, apply the same inequality to every selected--unselected pair;
their smallest initial gap is the rank-$k$ boundary gap.
\end{proof}

For affine scores $s_i(t)=a_i^Tt+b_i$ and any norm with dual norm
$\|\cdot\|_*$, the sharper sufficient condition
\begin{equation}
 s_r(t)-s_i(t)>\rho\|a_r-a_i\|_*
 \quad (i\ne r)
\end{equation}
certifies the same winner for every $\|t'-t\|\leq\rho$, by the dual-norm
inequality. At a tied winner no positive radius follows from these strict-gap
conditions. In a graph, all relevant node bounds must concern the actual
perturbed candidate scores, including changes propagated from earlier nodes.
A local observed margin is not by itself a radius in input or parameter
space. Dropout masks and normalization exponent selections are additional
trace state; extrema margins alone do not certify their stability.

Even if a tape remains sufficient for the original execution after live
weights change, it is not automatically the derivative of a fresh execution.
For the additive operators above the recorded reverse coefficients do not
depend on weight values, so the original selected reverse action remains
well defined. General local evaluators may instead require the original
parameter snapshot in $S$.

\subsection{Conditional information and representation accounting}
\label{app:information}

Fix available information $S=s$. Regard two admissible forward executions as
equivalent when their requested selected reverse actions agree for every
admissible cotangent. Let $J$ denote this equivalence class and let $M_s$ be
the number of attainable classes. Counting route labels can overcount these
classes: different labels can select identical coefficients, cancel after
composition, or be indistinguishable for the requested subset of gradients.

\begin{proposition}[Conditional action-coding bound]
\label{prop:action-coding}
A deterministic exact decoder that must answer every admissible cotangent
from $(c,S)$ needs at least $\lceil\log_2 M_s\rceil$ bits for a fixed-length
certificate at $S=s$, when $M_s$ is finite. A conditionally uniquely decodable
variable-length encoding has expected length at least $H(J\mid S)$ under
the stated distribution of executions.
\end{proposition}

\begin{proof}
Two inequivalent executions admit a cotangent on which the requested actions
differ. They cannot share a code when the decoder's other input $S=s$ is
identical. Hence at least $M_s$ distinct codes are necessary, proving the
fixed-length bound. Conditional unique decodability gives the Kraft
inequality for each $s$; the standard coding inequality yields expected
length at least $H(J\mid S=s)$. Average over $S$ for the conditional bound.
\end{proof}

For a single additive extremum with an omitted freely varying input, fixed
weights, and all $K$ candidates attainable and derivative-distinguishable,
the local bound is $\lceil\log_2K\rceil$. For $n$ outputs it becomes
$\lceil n\log_2K\rceil$ only if all $K^n$ joint action classes are attainable
and distinguishable under the same conditioning. Independent per-output
fixed-width packing uses $n\lceil\log_2K\rceil$ bits, which need not attain
that joint bound. A top-$k$ sum has at most $\binom{K}{k}$ selected-set actions
per output, not $K^k$ distinct ordered actions. Thus storing $k$ indices is a
simple implementation choice, not a proof of information optimality.

Crucially, if $S$ already contains the full original input and parameters of
a deterministic program, its route and selected reverse action are
determined by $S$, so $H(J\mid S)=0$. Merely prohibiting forward
recomputation does not produce a positive information-theoretic lower bound
in that setting; a nontrivial time--space claim requires a computational
model. Proposition~\ref{prop:action-coding} is an interface-conditional coding
statement, not a universal lower bound on training memory.

\paragraph{Current adapter payload.}
\label{app:payload}
For batch size $B$, length $T$, width $D$, $L$ layers with tap counts $K_l$,
vocabulary size $V$, and a top-$k$ sum head, define
$\ell(n)=\lceil\log_2n\rceil$, with $\ell(1)=0$. The current hard MorphSSM
adapter's declared logical payload is
\begin{align}
 N_{\mathrm{bits}}=BT\Big[&D\sum_{l=1}^{L}\big(\ell(K_l)+\ell(D)\big)
       +Vk\ell(D)\notag\\
       &+5(2L+1)+\ell(V)+Dn_{\mathrm{drop}}\Big],
 \label{eq:morphssm-logical-bits}
\end{align}
where $n_{\mathrm{drop}}$ counts the actually recorded input/layer dropout
masks. Each layer contributes convolution and projection winners and two
normalization exponents; the final norm contributes one exponent, the head
contributes $BT Vk$ indices, and token IDs contribute $BT$ codes. Exponents
are clamped to $[-8,8]$, so their 17-code alphabet uses five bits.

For fields with $n_f$ entries of width $b_f$, device payload bytes are
$\sum_f\lceil n_fb_f/8\rceil$, with at most seven padding bits per field.
The portable contiguous payload is $\lceil N_{\mathrm{bits}}/8\rceil$;
its directory and framing add bytes. These formulas exclude logits,
cotangents, parameter/optimizer state, Python metadata, and transient
packing, decoding, or reduction workspace. They are not allocator-peak
formulas. The vocabulary-head term can dominate, so there is no
architecture-independent memory advantage over checkpointing.

\subsection{Exclusions and the separate arithmetic contract}
\label{app:exclusions}

\paragraph{Dense trainable maps.}
For $y=Wx$, a fixed unchanged $W$ suffices for the state action $W^Tq$, but
the weight action is $qx^T$. With $W$ fixed, two distinct omitted inputs and
a suitable common nonzero $q$ give different weight actions. A route-free
record cannot represent both under this interface. Dyadic forward weights
do not eliminate this activation dependence for gradients of real-valued
master weights under an explicitly chosen straight-through rule. A
discrete sign-update recipe is a different contract, not an exact ordinary
weight-gradient identity.

\paragraph{Positive temperature and smooth mixtures.}
For $f_\tau(x)=\tau\log\sum_i\exp(x_i/\tau)$ with $\tau>0$, the selected
maximum index does not determine the derivative
$\exp(x_i/\tau)/\sum_j\exp(x_j/\tau)$. For example, $(0,-a)$ has the same
maximum for every $a>0$ but continuously changing derivatives. Small positive
temperature is not the hard operator. Likewise, an expert index does not
determine a dense expert's parameter gradient, nor does a hard gate certify
the value-dependent derivatives of a smooth attention mixture.

\paragraph{Normalization.}
Ordinary normalization has a value-dependent Jacobian. For the simple
$y=x/\sqrt{x^2+\varepsilon}$, the derivative
$\varepsilon/(x^2+\varepsilon)^{3/2}$ is not determined by the sign of $x$;
variance-normalized layers have the analogous continuous dependence.
The implemented tropical layer norm instead declares its selected exponent
$s$ constant in differentiation and uses
$y=2^{-s}(x-\mathrm{mean}(x))+\beta$. Its local reverse action is
\begin{equation}
 \bar x=2^{-s}\left(q-\frac{\mathbf{1}^Tq}{D}\mathbf{1}\right),
 \qquad \bar\beta=\sum_{\text{rows}}q.
\end{equation}
Thus the exponent is sufficient auxiliary state, but not a winner route.
Exponent quantization can make the forward map discontinuous at thresholds;
the stop-gradient rule is a declared program rule, not a general classical
derivative there. The mean's factor $1/D$ is dyadic only when $D$ is a power
of two. A trainable multiplicative gain would also require the normalized
activation for its gain gradient.

\paragraph{Whole training recipes.}
Gather, scatter, copy, sign, and addition implement the additive selected
reverse actions without multiplying two run-time numerical data values.
That property is separate from trace sufficiency: an unchanged dense matrix
may be reverse-available yet require products, and even a fully routed trunk
may use a conventional softmax loss or AdamW outside its boundary. A software
audit must inspect the executed forward, backward, loss, and update under
its versioned operation semantics. A physical zero-DSP or timing claim
additionally requires evidence for a particular synthesized design; neither
the route theorem nor a software pass supplies that evidence.

\subsection{Why small gradient errors can change an AdamW trajectory}
\label{app:adamw-sensitivity}

Route sufficiency is a statement about reverse dependencies, not numerical
conditioning of the optimizer. For standard AdamW
\citep{loshchilov2019adamw}, zero initial moments, learning rate $\eta>0$
and $\epsilon>0$, bias correction gives the scalar first-step adaptive term
\begin{equation}
 u(g)=-\eta\frac{g}{|g|+\epsilon}.
 \label{eq:adamw-first-step}
\end{equation}
Decoupled weight decay adds the same term $-\eta\lambda\theta$ to two
executions starting from the same parameter; it does not remove this
sensitivity. For $g\ne0$,
$|u'(g)|=\eta\epsilon/(|g|+\epsilon)^2\le\eta/\epsilon$; the same
bound follows at zero by continuity of the derivative. The mean value
theorem therefore gives
\begin{equation}
 |u(g+\delta)-u(g)|\le\frac{\eta}{\epsilon}|\delta|.
\end{equation}
At $\eta=10^{-3}$ and $\epsilon=10^{-8}$ the local amplification bound is
$10^5$. In particular, sign-opposed cancellation residuals of order
$\epsilon$ can produce updates of order $\eta$. This elementary
conditioning observation is not specific to tropical models or a new
optimizer result. Hard selection adds a second source of sensitivity:
after parameters move, crossing a score tie can change the selected
derivative discontinuously.

It is consequently important to distinguish three checks: the correctness
of VJPs at identical model state and supplied cotangent; the optimizer's
behavior on identical supplied gradients and moment state; and agreement
of independently evolving finite-precision training trajectories. Changing
$\epsilon$ defines a different AdamW configuration and lowers this
\emph{first-step} bound, but does not guarantee long-horizon route identity,
learning quality or equivalent conditioning after moments have accumulated.
